%% ============================================================================
%% Capítulo 8 — Conclusão
%%
%% Fontes: exclusivamente os Capítulos 2-7 desta dissertação (todos os
%% números citados aqui são retomadas de números já apresentados e
%% rastreados a artefatos nos capítulos de origem). Estrutura conforme
%% docs/blueprint-dissertacao.md (padrão Esdras: "Limitações e Ameaças à
%% Validade" na Conclusão) e docs/guia-escrita-didatica.md ("conselhos ao
%% praticante" em checklist, à la Hutter §14.1).
%% ============================================================================

\chapter{Conclusão}
\label{chap:conclusao}

Este capítulo responde às três perguntas de pesquisa com a evidência
dos Capítulos~\ref{chap:benchmark} a \ref{chap:destilacao}, enumera
as contribuições --- com os resultados negativos no mesmo estatuto dos
positivos ---, traduz os resultados em orientações ao praticante,
consolida as limitações e ameaças à validade do estudo como um todo e
indica trabalhos futuros.

\section{Considerações finais}
\label{sec:conc-finais}

Esta dissertação partiu da pergunta que organizou uma década de
literatura sobre dados tabulares (``árvores ainda vencem?'') e
terminou respondendo a uma pergunta diferente, porque a evidência
coletada mostrou que a original expirou. As três perguntas de pesquisa
enunciadas no Capítulo~\ref{chap:introducao} recebem aqui as respostas
diretas, cada uma com a evidência do capítulo que a respondeu.

\textbf{QP1} --- \emph{o empate de acurácia entre
as famílias de modelos tabulares sobrevive a um protocolo rigoroso e à
chegada da geração 2026 de foundation models?} O
Capítulo~\ref{chap:benchmark} responde em
dois tempos: dentro da geração anterior a 2026, sim --- o topo é
equivalência prática aferida sob os critérios declarados (8/10 equivalentes sob ROPE
$\pm0{,}01$ no critério brando $E[\theta]>0{,}5$, e 3/10 --- os
GBDTs --- no corte estrito de 0,95; sensibilidade
$0{,}005/0{,}01/0{,}02 \to 2/8/10$ no
critério brando), com a interação modelo$\times$conjunto
dominando a variância de \emph{rank} (0,69, contra 0,20 da família e
0,11 do modelo-dentro-da-família), de modo que o eixo ``GBDT vs.\
aprendizado profundo'' é uma abstração fraca para esses dados; entre
gerações, não --- o
paradigma pré-treinado separa-se do pelotão (TabFM \emph{zero-shot}
com \emph{ranks}
médios 2,7/1,0/1,2 nas três tarefas, melhor modelo em 13 dos 18
conjuntos; TabPFN como vice consistente),
sustentado pela consistência entre conjuntos e pelos tamanhos de
efeito pareados --- não pelo \emph{post-hoc} frequentista de baixo
poder nem pelo instrumento bayesiano, que não separa o TabFM de
XGBoost e CatBoost em equivalência prática --- e ao preço de quatro a
cinco ordens de magnitude em latência de serviço (43.262 contra
0,33~\us/linha do XGBoost no protocolo padronizado). A
camada por instância (Seção~\ref{sec:cap5-ato5}) diz ainda \emph{onde}
a quebra mora: XGBoost e \emph{foundation models} concordam em
${\ge}93{,}6\%$ das instâncias (mediana), a franja disputada é
convertida de forma fortemente heterogênea entre conjuntos, e as
explicações fáceis (margem baixa, classe minoritária, poucos dados)
não foram confirmadas sob os critérios pré-registrados.

\textbf{QP2} --- \emph{que artefato de decisão um benchmark com $N$
realista de conjuntos consegue de fato sustentar, um roteador
aprendido ou uma política por restrições --- e como essa resposta muda
quando a decisão pondera acurácia e custo?} Uma política por
restrições, não um roteador aprendido. Sob
\emph{leave-one-dataset-out} com 18 conjuntos, o roteador por
meta-características acerta menos que a regra trivial de classe
majoritária (0,11 contra 0,44 com 11 modelos; 0,72 contra 0,78 com
14), embora as reversões condicionais que ele deveria explorar sejam
reais --- o aprendizado profundo recente toma o topo
multiclasse (descritivo, $N{=}3$), os GBDTs ganham vantagem sob
desbalanceamento severo, o TabPFN domina o regime de poucos dados
($n \leq 4.000$; e a subamostragem intra-conjunto mostra que a
vantagem dos \emph{foundation models} não se restringe a conjuntos
pequenos, Seção~\ref{sec:cap5-ato5}). O que os dados sustentam é um
arcabouço de restrições
verificáveis antes de qualquer treino (latência, limites arquiteturais,
interpretabilidade), com a política ``\emph{foundation model} primeiro''
como \emph{default} de \emph{regret} mediano zero na métrica primária.
Quando a decisão pondera acurácia e custo, a resposta muda cedo: a
recomendação deixa de ser o \emph{foundation model} entre $w{=}0{,}9$
e $w{=}0{,}8$ de peso na acurácia, na normalização por postos adotada
(experimento E2, Capítulo~\ref{chap:framework}) --- e é por isso que a
primeira pergunta do fluxograma é a do orçamento de latência: é a que
mais cedo muda a resposta.

\textbf{QP3} --- \emph{o que a destilação de conhecimento transfere de
um foundation model tabular em regressão, a que custo de serviço, e sob
quais condições verificáveis a priori?} Transfere a forma da
distribuição preditiva, não a média. O essencial do que o professor
sabe sobre a média, um GBDT ajustado aprende sozinho dos rótulos
brutos: o aluno treinado nas médias do professor supera o controle em
4 de 20 conjuntos com o TabPFN, e a réplica com o TabFM como
professor, nos cinco conjuntos centrais, mantém o negativo. O
treinado nas curvas de quantis retém, na mediana, 19\% do gap de CRPS
com hiperparâmetros herdados (12\% com os dois braços re-ajustados),
com retenção positiva em 16 de 20 conjuntos (teste de sinal,
$p\approx0{,}019$ com a correção sequencial do desenho em duas
etapas), e supera o melhor entre QRF e NGBoost em 17 de 20
(experimento E1). O custo de serviço do aluno distribucional é da
ordem de $10^{2}$~\us{} por linha nos casos medidos da fronteira (99 e
169~\us/linha; Seção~\ref{sec:destil-fronteira}) --- acima do controle
pontual, e $2{,}5\times$ a $27\times$ abaixo da configuração
competitiva mais barata do professor onde a vantagem é grande. As
condições verificáveis são três, e a rotulagem \emph{out-of-fold} que
as mede custa minutos: transformar o alvo, testar um nativo
distribucional forte e destilar somente a vantagem que sobrar ao
professor --- com a própria rotulagem OOF como salvaguarda de
calibração do aluno (0,5--5,1~p.p.\ de PICP80,
Capítulo~\ref{chap:destilacao}).

\section{Contribuições}
\label{sec:conc-contribuicoes}

As cinco contribuições enunciadas na Seção~\ref{sec:intro-contribuicoes}
são retomadas aqui na mesma ordem; sob cada uma listam-se os
resultados que a compõem, com a natureza de cada um ---
\emph{evidência} (resultado empírico sob protocolo declarado),
\emph{recorte} (delimitação verificada de lacuna) ou \emph{regra}
(prescrição operacional derivada da evidência) --- e o capítulo de
origem. Os resultados negativos são listados com o mesmo estatuto dos
positivos, conforme o contrato de conduta do estudo
(Capítulo~\ref{chap:metodologia}).

\begin{enumerate}
  \item \textbf{O benchmark multicritério da geração 2026.}
  \begin{enumerate}
  \item \textbf{Benchmark neutro de 14 modelos
    (evidência; Capítulo~\ref{chap:benchmark}).} 14 modelos $\times$ 18
    conjuntos (250 das 252 células; duas exclusões estruturais), com
    \emph{hold-out} intocado pela seleção, inferência frequentista \emph{e}
    bayesiana com ROPE e análise de sensibilidade, e três eixos de custo
    sob o mesmo protocolo. Dele saem o empate positivo da geração
    anterior, o deslocamento do topo pela geração 2026 e, até onde alcançam as
    verificações
    datadas, a primeira comparação neutra
    TabPFN~$\times$~TabFM sob protocolo idêntico, enquadrada como
    replicação inter-laboratórios do paradigma, com versões fixadas e
    políticas espelhadas.
  \item \textbf{Análise de complementaridade por instância (evidência;
    Capítulo~\ref{chap:benchmark}, Seção~\ref{sec:cap5-ato5}).}
    Pré-registrada com critérios de falsificação datados e reprodução
    exata dos modelos finais do benchmark (52/52 células em \emph{gate}
    de sanidade), estabelece que XGBoost e \emph{foundation models}
    concordam em ${\ge}93{,}6\%$ das instâncias (mediana), que a
    vantagem do FM é fortemente heterogênea entre conjuntos (McNemar
    exato por par) e, como negativos informativos, não confirma as
    explicações de margem baixa, classe minoritária e vantagem restrita
    a poucos dados (esta última com sinal descritivo contrário no
    agregado).
  \end{enumerate}
  \item \textbf{O teste independente da família Kolmogorov--Arnold.}
  \begin{enumerate}
  \item \textbf{Resultado desfavorável sob protocolo uniforme
    (evidência negativa;
    Capítulo~\ref{chap:benchmark}).} No primeiro confronto da família
    com a grade completa GBDT/DL/\emph{foundation models} em protocolo
    multicritério --- até onde alcançam as verificações datadas ---, KAN
    e TabKAN terminam no último terço nas três tarefas
    (\emph{ranks} 11,0/12,0/9,8 e 11,4/13,3/11,8), em linha com a
    crítica de paridade de custo \citep{kancritical2024} e contra as
    alegações dos proponentes --- e contra o resultado favorável ao
    KAN de \citet{poeta2024}, obtido sem \emph{baselines} GBDT e sob
    treino curto, que o protocolo uniforme não reproduz; a contraparte registrada é que ambas são
    baratas de servir (2,3--3,1~\us/linha) --- o problema é desempenho,
    não custo.
  \end{enumerate}
  \item \textbf{O arcabouço de decisão avaliado.}
  \begin{enumerate}
  \item \textbf{Resultado negativo do roteamento por meta-características
    (evidência negativa; Capítulo~\ref{chap:framework}).} Sob LODO com
    $N{=}18$, um roteador aprendido não supera o \emph{baseline} de
    classe majoritária em nenhuma geração (0,11 vs.\ 0,44 com 11
    modelos; 0,72 vs.\ 0,78 com 14) --- corroborado por evidência
    independente contemporânea \citep{herre2026}. A consequência de
    desenho: o arcabouço não contém preditor de família vencedora.
  \item \textbf{Arcabouço de decisão avaliado por \emph{regret}
    (evidência~$+$~regra; Capítulo~\ref{chap:framework}).} A política
    fixa ``\emph{foundation model} primeiro'' é regret-ótima na métrica
    primária (\emph{regret} mediano 0,000 na geração 2026,
    preservado na política executável sempre-TabFM --- média 0,080,
    pior caso 1,000 no \texttt{telco}, este último em grande parte
    artefato do pré-processamento medido no
    Capítulo~\ref{chap:metodologia}), e o
    arcabouço final é um fluxograma de quatro restrições verificáveis
    \emph{a priori} --- a forma que os dados sustentam.
  \item \textbf{Regret multicritério ponderado --- E2 (evidência;
    Capítulo~\ref{chap:framework}).} O
    experimento registrado que localiza o ponto de inversão da
    recomendação: ``FM primeiro'' só é ótima na mediana para
    $w \geq 0{,}9$ de peso na acurácia; com $w \leq 0{,}6$,
    ``sempre-GBDT'' tem \emph{regret} mediano zero. É a justificativa
    quantitativa da arquitetura do fluxograma --- a latência é a
    primeira pergunta porque é a que mais cedo muda a resposta.
  \end{enumerate}
  \item \textbf{A destilação distribucional de \emph{foundation
    models} tabulares.}
  \begin{enumerate}
  \item \textbf{O recorte e o resultado central (recorte~$+$~evidência;
    Capítulo~\ref{chap:destilacao}).} O recorte --- transferir curvas de
    quantis \emph{de um professor foundation model tabular} para um
    aluno GBDT, avaliando por CRPS e cobertura --- foi verificado como
    sem trabalho \emph{acadêmico} anterior nas buscas datadas deste
    estudo, de julho a novembro de 2026 (registro completo no
    Apêndice~\ref{app:suplementares}; um
    repositório público de setembro de 2026 --- referenciado na
    Seção~\ref{sec:rel-destilacao} --- destila o TabFM em XGBoost
    para classe e regressão \emph{pontual}, sem tocar o recorte
    distribucional --- a linhagem não tabular de destilação
    distribucional é citada como antecedente na
    Seção~\ref{sec:rel-destilacao}). A evidência: retenção mediana de
    12\% a 19\% do gap de CRPS, conforme o ajuste dos braços,
    positiva em 16 de 20 conjuntos ($p=0{,}006$, $p\approx0{,}019$ sob
    a correção sequencial; probabilidade posterior
    bayesiana de 0,96 em retenção positiva, com massa posterior
    esperada 0,70), incluindo um aluno
    \emph{born-again} (que supera o próprio professor), servindo na
    casa de cem microssegundos por linha.
  \item \textbf{Refutação da destilação pontual em \emph{pool} realista
    (evidência negativa; Capítulo~\ref{chap:destilacao}).} O aluno
    pontual destilado supera o controle em apenas 4 de 20 conjuntos
    com o TabPFN; a réplica com o TabFM --- cuja vantagem
    pontual sobre o \emph{baseline} é quase $6\times$ maior no
    conjunto medido (o \texttt{california\_housing}, de alvo
    binarizado) --- mantém o negativo nos cinco conjuntos centrais, e a
    varredura de tamanho de \emph{pool} isola o
    moderador (ganho residual restrito a $n \leq 800$, dependente do
    conjunto) --- uma hipótese de reconciliação com os ganhos pontuais
    contemporâneos de \citet{jeong2026}, cuja receita difere da nossa
    em eixos que não replicamos (rotulagem com contexto completo,
    consultas sintéticas, aluno treinado só em predições do
    professor).
  \item \textbf{Fronteira de três estratégias de \emph{serving}
    (evidência; Capítulo~\ref{chap:destilacao}).} Destilar, comprimir
    contexto e reduzir \emph{ensemble} medidos pela primeira vez nos
    mesmos \emph{hold-outs}: onde a vantagem do professor é grande, o
    aluno destilado é, entre as estratégias derivadas do professor, o
    sistema distribucional CRPS-ótimo abaixo da
    configuração competitiva mais barata dele ($2{,}5\times$ e
    $27\times$ a latência do aluno nos dois casos medidos --- o E1
    mostra que uma QRF nativa também habita o regime); onde é
    pequena, reduzir o \emph{ensemble} basta. Dois achados colaterais de
    engenharia: um membro do TabPFN $\approx$ oito a um sétimo da
    latência, e a última duplicação de contexto é desproporcionalmente
    cara em GPUs de 16~GB.
  \item \textbf{Ablação desconfundida da rotulagem \emph{out-of-fold}
    (evidência e regra; Capítulo~\ref{chap:destilacao}).} Em regressão, o
    vazamento do professor não degrada acurácia --- degrada
    \emph{calibração} (perda de 0,5 a 5,1~p.p.\ de PICP80 a contexto
    idêntico): um refinamento desconfundido do achado de classificação
    do Pocket~FM \citep{pocketfm2026}.
  \item \textbf{Confronto com nativos distribucionais clássicos --- E1
    (evidência; Capítulo~\ref{chap:destilacao}).} O aluno destilado vence o
    melhor entre QRF e NGBoost em 17
    de 20 conjuntos ($p=0{,}0013$), com os três reveses diagnosticáveis
    --- no \texttt{wine\_quality}, onde o \emph{gate} pré-registrado já
    desaconselhava destilar, a QRF \emph{default} supera até o
    professor. O complemento negativo tem valor próprio: o passo ``teste
    um nativo forte antes'' da regra de decisão deixou de ser retórico.
  \item \textbf{Negativo arquitetural do TabFM como professor
    distribucional (evidência negativa;
    Capítulo~\ref{chap:destilacao}).} A cabeça de regressão do TabFM é
    pontual --- verificação arquitetural de julho de 2026, confirmada
    pelo relatório técnico oficial \citep{kong2026} --- e a dispersão de
    seu \emph{ensemble} é mensuravelmente inútil como distribuição
    preditiva (CRPS 0,524 contra 0,069 do TabPFN; cobertura de 16,3\% a
    nominal de 80\%): líder pontual não é líder distribucional, e a
    cabeça distribucional é o ativo arquitetural decisivo para a
    destilação.
  \item \textbf{Regra de decisão para destilar (regra;
    Capítulo~\ref{chap:destilacao}).} Derivada de duas hipóteses que
    não se sustentaram (ponto: refutada; cauda pesada: não confirmada)
    e do E1: transformar o alvo, testar
    um nativo forte, e destilar somente a vantagem genuína que
    sobreviver --- pré-condição mensurável por uma rotulagem OOF de
    minutos. É a forma operacional da resposta à QP3.
  \end{enumerate}
  \item \textbf{Infraestrutura e conduta reutilizáveis (transversal;
    Capítulos~\ref{chap:metodologia} a~\ref{chap:destilacao}).} Um
    protocolo versionado de ponta a ponta --- configurações, sementes,
    \emph{hashes} de dados, pré-registros com adendos datados,
    verificações de novidade datadas e errata pública --- em que todo
    número do texto rastreia a um artefato. A contribuição aqui é de
    \emph{processo}: o contrato de conduta determinou decisões concretas
    de escopo, e seus mecanismos de correção foram exercitados em casos
    reais.
\end{enumerate}

\section{Conselhos ao praticante}
\label{sec:conc-conselhos}

O destilado operacional dos Capítulos~\ref{chap:framework}
e~\ref{chap:destilacao}, na ordem em que as perguntas devem ser feitas
--- cada item rastreia ao resultado que o sustenta.

\begin{itemize}
  \item \textbf{Comece pelas restrições, não pelo desempenho.} Quatro
    perguntas verificáveis antes de qualquer treino: orçamento de
    latência, limites arquiteturais, exigência de interpretabilidade,
    e só então o \emph{default} (fluxograma da
    Seção~\ref{sec:arcabouco-resultante}). Não tente aprender um
    roteador por meta-características com poucas dezenas de conjuntos: sob
    LODO, ele não supera a regra trivial
    (Seção~\ref{sec:resultado-negativo}).
  \item \textbf{Se latência de microssegundos importa, GBDT.} Com 40\%
    ou mais do peso da utilidade na latência, ``sempre-GBDT'' tem \emph{regret}
    mediano zero; a recomendação ``FM primeiro'' já se inverte com
    10--20\% de peso na latência (E2,
    Seção~\ref{sec:regret-ponderado}).
  \item \textbf{Se nenhuma restrição está ativa, \emph{foundation model}
    primeiro.} Regret mediano 0,000 na geração 2026
    (Seção~\ref{sec:resultado-positivo}) --- com o TabPFN como opção
    madura e o TabFM como fronteira, nas versões e com as ressalvas de
    maturidade declaradas no Capítulo~\ref{chap:benchmark}.
  \item \textbf{Verifique os limites arquiteturais cedo.} Mais de 10
    classes exclui o TabFM; 100 classes excluem ambos os
    \emph{foundation models} (o caso \texttt{helena}); estouro de
    contexto é restrição objetiva --- são exclusões estruturais, não
    detalhes (Seção~\ref{sec:arcabouco-resultante}).
  \item \textbf{Em regressão com incerteza: transforme antes de
    destilar.} Em alvos de escala de preço, uma transformação
    logarítmica de uma linha é gratuita e muda o regime de erro de
    todos os sistemas --- ela não elimina a vantagem relativa do
    professor, mas é o primeiro teste a fazer, e custa minutos
    (Seção~\ref{sec:destil-regra}).
  \item \textbf{Teste a QRF antes de destilar.} Um nativo distribucional
    \emph{default} custa segundos de treino, serve a microssegundos e,
    numa minoria previsível de conjuntos, resolve o problema antes que
    qualquer professor precise ser carregado --- no conjunto
    \texttt{wine\_quality}, venceu até o professor (E1,
    Seção~\ref{sec:destil-e1}); um MLP-quantil cobre os conjuntos
    pequenos e numéricos (Seção~\ref{sec:destil-regra}).
  \item \textbf{Destile quando a vantagem sobreviver às
    transformações e a um nativo forte --- e meça-a antes.} Uma única
    rotulagem OOF, de custo de minutos, mede diretamente a vantagem
    genuína do professor; destile somente se ela persistir
    (Seção~\ref{sec:destil-regra}). Não destile a média: em pools
    realistas, não há o que transferir no eixo pontual
    (Seção~\ref{sec:destil-ponto}).
  \item \textbf{Rotule sempre \emph{out-of-fold}.} O custo é modesto, e
    a salvaguarda protege o que as métricas de acurácia não veem: a
    cobertura dos intervalos do aluno (0,5--5,1~p.p.\ de PICP80;
    Seção~\ref{sec:destil-oof}).
  \item \textbf{Antes de pagar o professor completo, reduza o
    \emph{ensemble}.} Um único membro do TabPFN aproximou os oito a um
    sétimo da latência --- às vezes sem custo em acurácia --- e, quando a
    vantagem
    do professor é pequena, essa alavanca encerra o caso sem destilação
    alguma (Seção~\ref{sec:destil-fronteira}).
  \item \textbf{Em GPUs de 16~GB, desconfie do último dobro de
    contexto.} A duplicação final, de 25~mil para 50~mil linhas, custou
    $3{,}3\times$ a latência por ganho marginal de acurácia
    (Seção~\ref{sec:destil-fronteira}); reserve o professor de contexto
    completo para processamento em lote.
\end{itemize}

A lista acima ordena as perguntas; falta nomear os cenários. A
Tabela~\ref{tab:conc-perfis} consolida os mesmos resultados como
\emph{perfis de uso}: para cada condição verificável antes de qualquer
treino, a recomendação que os dados sustentam, a evidência medida que a
sustenta e --- a coluna que os \emph{benchmarks} usuais omitem --- o que
se aceita perder em troca. A coluna final é o ponto: em vários perfis, a
recomendação é deliberadamente um modelo \emph{pior na métrica
primária}, e a troca é defensável porque foi medida. Dois resultados
sustentam essa contabilidade: no pelotão do núcleo, abrir mão do topo
custa tipicamente menos de um ponto percentual de AUC
(Seção~\ref{sec:cap5-ato1}), ao passo que as diferenças de latência,
de \emph{recall} da classe minoritária e de custo de cauda são de
fatores inteiros a ordens de magnitude (colunas 3 e 4 da tabela).
Quando a métrica em que os modelos diferem é a que a
aplicação paga, é ela que decide --- e o quadro diz qual é, caso a
caso.

\begin{table}[htbp]
\centering
\caption[Perfis de uso: condição, recomendação, evidência e troca]{Perfis de uso: condição verificável, recomendação, evidência e
custo aceito em troca. Fonte: Capítulos~\ref{chap:benchmark}--\ref{chap:destilacao}, seções indicadas.}
\label{tab:conc-perfis}
\scriptsize
\setlength{\tabcolsep}{5pt}
\renewcommand{\arraystretch}{1.15}
\begin{tabular}{p{0.21\linewidth}p{0.17\linewidth}p{0.33\linewidth}p{0.19\linewidth}}
\toprule
\textbf{Perfil (condição)} & \textbf{Recomendação} & \textbf{Evidência} &
\textbf{Aceita-se em troca} \\
\midrule
Nenhuma restrição ativa; latência tolerante &
\emph{Foundation model} primeiro (TabPFN maduro; TabFM fronteira) &
\emph{Regret} mediano 0,000 com 14 modelos na avaliação do
Capítulo~\ref{chap:framework} (teto da família, medido nos mesmos 18
conjuntos; a
executável sempre-TabFM preserva a mediana, com média 0,080 e pior
caso 1,0; sempre-TabPFN: mediana 0,242); melhor absoluto em 13 de 18
(Seções~\ref{sec:resultado-positivo}, \ref{sec:cap5-ato4}) &
Custo de serviço (de ${\sim}3.900\times$ o GBDT mais lento ---
TabPFN sobre LightGBM, Tabela~\ref{tab:app-custos} --- a
$131.097\times$ o XGBoost); licença não comercial
do TabPFN-2.5 \\
\addlinespace
Poucos dados ($n \lesssim 4.000$) &
TabPFN &
Rank médio 1,0 no regime, contra 2,8 do XGBoost, nos 5 conjuntos
varridos (Seção~\ref{sec:cap5-ato3}) &
Latência de inferência; a liderança com $n$ grande depende do conjunto \\
\addlinespace
Classe rara com falso negativo caro &
FM a limiar fixo --- \emph{ou} re-otimizar o limiar do GBDT; medir
recall a precisão igualada antes de trocar &
\texttt{amazon\_employee}: acurácias 0,943--0,949 indistinguíveis;
recall da minoritária 0,032 $\to$ 0,298 no limiar 0,5
(Seção~\ref{sec:cap5-ato5}). Contrapesos: a hipótese de ganho
concentrado na classe rara não se confirmou; o
melhor AUC do conjunto é de um GBDT (CatBoost 0,887) --- vantagem de
limiar fixo, em um conjunto &
Precisão (0,706 $\to$ 0,621) e latência; a vantagem é frequente, não
lei (\texttt{give\_me\_some\_credit}: recalls a $<$0,02) \\
\addlinespace
Regressão com decisão de cauda e orçamento de \us &
Aluno destilado (correção conformal da cobertura: piloto em 5
conjuntos) &
Retenção positiva em 16/20 ($p{=}0{,}006$; $0{,}019$ sequencial),
mantida re-tunada; em $\tau{=}0{,}90$, 18/20 ($p{=}0{,}0002$), ganho
${\sim}8\%$ nas caudas; combinação destilado$+$conformal não medida
junta nos 20 (Seções~\ref{sec:destil-dist},
\ref{sec:destil-e3}) &
Nada no eixo pontual (4 de 20); latência acima do controle, na casa
de $10^2$~\us{} por linha; \emph{pipeline} OOF; a licença não
comercial do professor alcança o aluno destilado \\
\addlinespace
A métrica da aplicação inverte o vencedor &
Decidir pela métrica que a aplicação paga; reportar as duas camadas &
\texttt{telco}: TabFM perde acurácia (0,771 vs.\ 0,801) e a franja no
McNemar ($-29{,}8$/mil), mas vence recall (0,663 vs.\ 0,570) e
balanceada (0,737 vs.\ 0,727) (Seção~\ref{sec:cap5-ato5}) &
O ``vencedor'' da primária; limiar não otimizado; e o
\texttt{TotalCharges} (cap.~4) alcança este conjunto: com a coluna
numérica, o TabFM salta de último a melhor AUC entre os 14
reavaliados \\
\addlinespace
Vantagem heterogênea por conjunto (franja de discordância por
instância, Seção~\ref{sec:cap5-ato5}) &
Escolha por conjunto/tarefa, nunca por instância; sem roteador
aprendido &
Concordância $\geq$0,936; TabFM converte a franja no \texttt{jannis}
($+59{,}6$/mil) e a perde no \texttt{telco}; roteador LODO abaixo do
trivial (Seções~\ref{sec:cap5-ato5}, \ref{sec:resultado-negativo}) &
A promessa de ``um vencedor universal''; a vantagem não encolhe com o
tamanho \\
\addlinespace
Latência de microssegundos (serviço \emph{online}) &
GBDT nativo; destilação recupera parte da vantagem distribucional &
0,33~\us/linha (XGBoost) vs.\ 43.262 (TabFM), $131.097\times$;
recomendação FM inverte com 10--20\% de peso na latência (E2)
(Seções~\ref{sec:cap5-ato2}, \ref{sec:regret-ponderado}) &
Rank do topo --- perda $<$1~p.p.\ de AUC no pelotão do núcleo (ROPE
vs.\ TabPFN: 8/10 no critério brando, 3/10 no estrito); vs.\ TabFM a
equivalência é
incerta ($P$ 0,33; $E[\theta]$ 0,40) \\
\addlinespace
Interpretabilidade intrínseca exigida (regulação) &
GBDT &
Única família forte em latência \emph{e} interpretabilidade na matriz
multicritério --- atribuição qualitativa por família, declarada, não
medida; portão 3 do fluxograma
(Seção~\ref{sec:arcabouco-resultante}) &
Rank do topo, sob a mesma contabilidade ROPE; sem destilação (o aluno
herda rótulos de um professor opaco) \\
\addlinespace
Limite arquitetural ($>$10 classes; contexto estourado) &
GBDT ou DL recente conforme a tarefa &
\texttt{helena} (100 classes) exclui os dois FMs; topo multiclasse
descritivo: STab/TabM 3,67 vs.\ XGBoost 6,0, $N{=}3$
(Seções~\ref{sec:cap5-ato3}, \ref{sec:arcabouco-resultante}) &
Nada: a exclusão é estrutural, não preferência \\
\addlinespace
Vantagem pequena do professor; alvo de preço; nativo forte disponível &
\emph{Não} destilar: QRF primeiro, $\log(1{+}y)$, redução de
\emph{ensemble} &
E1: QRF vence o destilado em 3 de 20 (no \texttt{wine\_quality},
até o professor); \texttt{kin8nm}: 1 membro à latência do aluno; o log
muda o regime de erro de graça, sem eliminar a vantagem relativa do
professor (Seções~\ref{sec:destil-e1},
\ref{sec:destil-fronteira}, \ref{sec:destil-regra}) &
O \emph{pipeline} de destilação inteiro --- e é ganho, não perda \\
\bottomrule
\end{tabular}
\end{table}

Três leituras fecham o quadro. Primeiro, nos perfis em que a
recomendação contraria a métrica primária (latência,
interpretabilidade, classe rara), a troca está contabilizada nas duas
colunas finais: a perda na primária
é limitada por medição (pelo ROPE no pelotão, pelo limiar declarado
nas operacionais), enquanto o ganho na métrica da aplicação é de
fatores inteiros a ordens de magnitude. Segundo, os perfis negativos
são tão informativos quanto os positivos: o quadro diz quando
\emph{não} usar o \emph{foundation model} (microssegundos, 100
classes, regulação) e quando \emph{não} destilar (\texttt{kin8nm},
\texttt{wine\_\allowbreak quality}), com os casos medidos no próprio
\emph{benchmark}. Terceiro, nenhuma linha é um
\emph{ranking}: seis são condições que se conhecem antes de qualquer
treino, na forma que o Capítulo~\ref{chap:framework} avaliou; as
outras quatro (classe rara, métrica que inverte, franja por instância,
vantagem do professor) pedem uma medição barata e declarada --- uma
segunda métrica no \emph{hold-out}, uma rotulagem
\emph{out-of-fold} --- antes de decidir.

\section{Limitações e ameaças à validade}
\label{sec:conc-limitacoes}

Cada capítulo declarou suas limitações locais; esta seção consolida as
do estudo como um todo, organizadas pelos três tipos canônicos de
validade (Seção~\ref{sec:met-limitacoes}). Duas limitações
transversais entram na taxonomia com destaque próprio: a
multiplicidade e a subcobertura do aluno destilado puro, abaixo, em
validade interna; a licença dos professores, em validade externa.

\subsection{Validade interna}

\textbf{Priors cientes de benchmarks --- o risco residual não
endereçável.} O TabFM declara pré-treino exclusivamente sintético; o
\emph{checkpoint} padrão do classificador TabPFN v2.5, porém, é
afinado em 43 conjuntos reais que incluem o \texttt{higgs} deste
benchmark (com a
sensibilidade medida nos doze conjuntos de classificação: diferenças
de AUC entre $-0{,}0051$ e $+0{,}0002$, e no \texttt{higgs} o
checkpoint só-sintético é até ligeiramente melhor, 0,8113 contra
0,8111; Seção~\ref{sec:met-limitacoes}); o regressor padrão é sintético puro,
preservando a parte distribucional. Permanece, além disso, o risco mais fraco de
\emph{priors} e políticas de inferência iterativamente ajustados, ao
longo das versões, ao desempenho em benchmarks públicos: uma forma
branda de sobreajuste da comunidade da qual nenhum benchmark público
escapa. O tratamento foi declarar o risco, enunciar conclusões no nível
de regime (consistência entre conjuntos, não células individuais) e
manter o eixo de custo, imune a contaminação, em toda conclusão
prática (Seção~\ref{sec:met-limitacoes}).

\textbf{Assimetrias declaradas nos experimentos complementares.} No E1,
NGBoost e QRF rodam em configuração \emph{default} documentada, sem
ajuste específico, e a QRF recebe imputação por mediana (ajustada no
\emph{pool}, aplicada ao teste) nos três conjuntos com valores ausentes; a
interpretação autorizada é ``\emph{default} forte'', não ``clássico
ajustado ao limite'' (Seção~\ref{sec:destil-e1}). Nos alunos de
destilação, a capacidade é herdada do \emph{tuning} do benchmark, sem
re-ajuste por alvo de destilação --- um controle deliberado, com a
ressalva declarada de que entre os braços mudam o rótulo \emph{e} a
perda; o custo da assimetria foi medido pelo
re-ajuste pré-registrado nos 20 conjuntos: o
veredito agregado sobrevive ao controle re-ajustado (16 de 20, mediana
0,124 contra 0,192 herdada, queda declarada), restando a retenção por
conjunto como dependente do orçamento de ajuste
(Seção~\ref{sec:destil-retune}).

\textbf{Erros de identidade de conjuntos.} Dois erros de rotulagem
herdados da criação do projeto (\texttt{diamonds}/\texttt{kin8nm};
``superconduct''/\texttt{online\_news}) foram descobertos durante o
estudo e corrigidos com re-derivação de todos os números afetados;
todo número desta dissertação é posterior às correções
(Seção~\ref{sec:met-datasets}; registro no
Apêndice~\ref{app:suplementares}). A ameaça residual é a que
correções não removem: a possibilidade de erros análogos ainda não
descobertos; a mitigação é estrutural (\emph{hashes} de dados,
configurações versionadas, números rastreáveis a artefatos).

\textbf{Multiplicidade entre experimentos.} Nenhuma correção é
aplicada \emph{entre} experimentos; sob correção de Holm na família
confirmatória, as conclusões primárias sobrevivem e os instrumentos
suplementares não --- e a manchete da QP3 carrega a correção
sequencial do seu desenho em duas etapas ($p\approx0{,}019$;
Seção~\ref{sec:destil-desenho}).

\textbf{Subcobertura do aluno destilado puro.} A cláusula
pré-registrada de cobertura (PICP80 a $\pm5$~p.p.\ do nominal) é
cumprida em só 10 dos 20 conjuntos --- a resposta à QP3 é
qualificada por isso: o ganho de calibração é comparativo, e o
fechamento da cobertura vem do passo conformal, pilotado em 5
conjuntos (Seção~\ref{sec:destil-dist}).

\subsection{Validade externa}

\textbf{$N$ pequeno por tarefa e os pisos aritméticos.} Com 10/3/5
conjuntos por tarefa, o poder dos testes é baixo por construção: em
regressão, o menor p-valor bilateral do Wilcoxon é 0,0625
(significância a 5\% é matematicamente impossível, antes de qualquer
correção) e a multiclasse ($N{=}3$) é tratada como descritiva em todo
o texto. As mitigações declaradas (equivalências sustentadas pela
camada bayesiana, nenhuma conclusão apoiada em significância onde não há
poder) limitam o dano, mas não o eliminam: as conclusões por tarefa
têm a precisão que $N$ permite, não mais
(Seção~\ref{sec:met-limitacoes}).

\textbf{A exclusão dupla do \texttt{helena}.} Os \emph{ranks}
multiclasse dos dois \emph{foundation models} são médias sobre apenas 2
conjuntos, fato repetido em toda tabela afetada. A alternativa
(remover o conjunto) esconderia exatamente o tipo de restrição
arquitetural que um estudo de decisão prática deve expor; o preço é que
a liderança multiclasse da geração 2026 repousa sobre a menor base
amostral de todo o benchmark.

\textbf{TabFM em versão fixada, com relatório técnico não revisado.}
Todos os resultados referem-se à versão v1.0.1 (julho de 2026, commit
fixado); o relatório técnico oficial \citep{kong2026} apareceu no fim de
setembro de 2026 e, até a escrita desta dissertação, não passou por
revisão por pares \citep{tabfm2026}. As conclusões sobre o modelo são
sobre \emph{essa} versão; e a sensibilidade de versão é, ela própria,
um trabalho futuro (Seção~\ref{sec:conc-futuros}).

\textbf{Regime delimitado.} Dezoito conjuntos OpenML (mais os quinze da
extensão CTR23 na destilação), teto de 100~mil amostras e uma única GPU
de 16~GB delimitam o regime estudado a dados tabulares de pequeno e
médio porte em hardware de porte comum. As ordens de grandeza entre
famílias tendem a transferir; os valores absolutos, não.

\textbf{O campo move-se em ciclos de meses.} TabPFN-2.5 em novembro de
2025, Pocket~FM em maio de 2026, TabFM em junho, a colisão parcial de
\citet{jeong2026} em outubro; a diferença entre as colunas @11 e @14
da validação LODO é a demonstração interna de quão rápido as conclusões
datam. Por isso toda afirmação de novidade e toda caracterização de
lacuna carrega data de verificação, nas buscas datadas deste estudo,
de julho a novembro de 2026 (registro completo no
Apêndice~\ref{app:suplementares}), e deve ser lida como ``aberta até
onde verificamos, nas datas declaradas'', nunca como certeza atemporal
(Capítulo~\ref{chap:related}).

\textbf{O alvo binarizado no núcleo de regressão.} O
\texttt{california\_\allowbreak housing} do núcleo é a variante de
alvo binarizado do \emph{tabular-benchmark} --- degenerada para
afirmações distribucionais. Cada resultado-manchete da destilação
carrega análise de sensibilidade com o conjunto excluído: o teste de
sinal, o E3, o Friedman de cinco sistemas e as ablações sobrevivem; o
Wilcoxon unilateral da retenção passa de $0{,}045$ a $0{,}067$ e deixa
de rejeitar a 5\%, degradação declarada com a mesma proeminência do
resultado (Capítulo~\ref{chap:destilacao}).

\textbf{Licença dos professores.} Os pesos do TabPFN-2.5 e do TabFM
são não comerciais, e a licença do TabPFN veda ainda destilar das suas
saídas um modelo ``competitivo''; o aluno destilado deste estudo opera
em regime de pesquisa, e implantação comercial exigiria licenciamento
próprio (Seção~\ref{sec:destil-limitacoes};
Apêndice~\ref{app:suplementares}).

\subsection{Validade de construto}

\textbf{O \emph{proxy} de latência do \texttt{adult}.} A latência entra
nas análises (inclusive no score ponderado do E2) pela medição
padronizada em um único regime (lote amortizado, cada modelo em seu
dispositivo natural), com o \texttt{adult} como referência. A direção do
viés foi declarada no pré-registro e a sua magnitude foi depois medida
nos extremos de tamanho do benchmark: os \emph{rankings} entre os 14
modelos são estáveis entre conjuntos (Spearman 0{,}78--0{,}95) e o
cruzamento do E2 é idêntico sob os três \emph{proxies}
(Seções~\ref{sec:met-limitacoes} e~\ref{sec:regret-ponderado}). O que
permanece não medido é a sensibilidade ao hardware: todos os números vêm
de uma única máquina declarada.

\textbf{O estimador de CRPS em grade.} A grade de 19 quantis
sobre-estima o CRPS exato em ${\sim}$4--5\% quase constantes, validado
contra a forma fechada gaussiana: comparações internas e razões de
retenção ficam preservadas, mas os valores absolutos de CRPS desta
dissertação não devem ser comparados aos de outros trabalhos
(Seção~\ref{sec:destil-metricas}).

\textbf{Métricas comprimem utilidades.} A métrica primária
\emph{threshold-free} e o score do E2 (dois critérios, pesos lineares)
são construtos deliberadamente simples de funções de utilidade que, na
prática, incluem custo de treino, robustez, interpretabilidade e
restrições regulatórias com pesos não lineares. O arcabouço mantém os
eixos separados em vez de agregá-los prematuramente, mas qualquer
recomendação escalar herda essa compressão
(Seções~\ref{sec:met-limitacoes} e~\ref{sec:limitacoes-framework}).

\section{Trabalhos futuros}
\label{sec:conc-futuros}

Listam-se apenas direções com lastro direto na evidência dos capítulos
(extensões que este estudo deixa armadas, não especulações).

\begin{itemize}
  \item \textbf{Professores mais novos: sensibilidade de versão da
    destilação.} O estudo fixou o TabPFN v2.5 por comparabilidade com o
    benchmark; TabPFN-3 e 3.5 \citep{tabpfn3-2026} já existem, e a
    colisão parcial de \citet{jeong2026} (ganhos pontuais com
    professores mais novos, sob protocolo distinto) sugere que a
    fronteira professor-aluno é sensível à versão. Repetir o protocolo
    distribucional com os professores da geração seguinte mediria
    diretamente se a retenção de 12--19\% cresce com a força do professor
    --- e se o negativo pontual resiste.
  \item \textbf{Calibração conformal como complemento e competidor.} A
    rota conformal pós-hoc sobre o próprio professor \citep{park2026}
    entrega garantias de cobertura sem aluno rápido --- ao preço da
    latência do \emph{foundation model}. Incorporá-la à fronteira (e
    aplicá-la também sobre o aluno destilado) poria à prova a vantagem
    de calibração observada na
    Seção~\ref{sec:destil-oof} contra um método com garantia formal.
  \item \textbf{Quantização como quarto braço da fronteira.} A
    quantização da atenção \citep{kubler2026} é a alavanca de
    \emph{serving} emergente que não medimos; a fronteira de três
    estratégias da Seção~\ref{sec:destil-fronteira} está construída
    para recebê-la como quarto braço nos mesmos \emph{hold-outs}.
  \item \textbf{Regret multicritério com preferências elicitadas.} O E2
    varreu pesos lineares sobre dois critérios, e a robustez do
    \emph{proxy} de latência ao conjunto de medição já foi verificada
    (Seção~\ref{sec:regret-ponderado}); a extensão natural é incorporar
    custo de treino, robustez e interpretabilidade ao score e elicitar
    funções de utilidade de praticantes reais --- testando se o ponto de
    inversão (10--20\% de peso na latência) se move com critérios
    adicionais.
  \item \textbf{KANs sob ajuste dedicado.} O resultado negativo do
    Capítulo~\ref{chap:benchmark} usa o protocolo uniforme da grade ---
    a leitura justa para a pergunta do benchmark, mas não a última
    palavra sobre a família: um estudo dedicado, com orçamento de busca
    próprio, o otimizador original e variantes de base, separaria
    definitivamente método de receita --- e diria se o último terço é
    teto ou artefato de protocolo.
  \item \textbf{Atualização viva do benchmark.} A infraestrutura
    resumível, com versões e sementes fixadas, está pronta para
    re-execução periódica à medida que a geração 2027 chegar; a
    integração com um benchmark vivo mantido pela comunidade, ao estilo
    TabArena \citep{tabarena2025}, é o caminho para que a profundidade
    de protocolo deste estudo sobreviva à sua datação --- a principal
    ameaça externa declarada acima.
\end{itemize}

\vspace{1em}

Em síntese, a dissertação entrega três peças: um benchmark datado da
geração 2026, com os eixos de custo medidos nas mesmas condições que a
acurácia; um arcabouço de restrições avaliado fora da amostra, com o
ponto de inversão por custo localizado; e uma regra, testada também
nos seus fracassos, para decidir quando destilar a distribuição
preditiva de um \emph{foundation model} em um aluno que serve na
casa de cem microssegundos por linha. As três dependem das versões de modelo, dos conjuntos
e do hardware declarados, e devem ser reavaliadas quando eles mudarem;
o protocolo público, com sementes, artefatos e erratas versionados,
existe para que essa reavaliação, e qualquer correção que ela traga,
custe pouco a quem vier depois.
